\chapter{Evidence of Required Ethics Training}

[Insert the student's evidence of completing all required ethics training.]

\chapter{Confirmation of Ethical Approval Waiver}

[Insert the confirmation email showing that ethical approval was waived.]

\chapter{Supplementary Experimental Results}

\section{Clothing-range coverage in the retained FITB subset}

The retained official FITB subset contained 46 menswear-labelled questions and
4,422 womenswear-labelled questions. Descriptive accuracy was 0.4348 for the
menswear subset and 0.6043 for the womenswear subset. The menswear result is not
treated as a reliable group comparison because its sample is substantially
smaller.

\chapter{Supplementary Software-Design Artefacts}
\label{app:software-design}

This appendix documents the implemented prototype at a level below the
methodology diagrams. Cove does not contain accounts, administrator roles,
repository interfaces or relational database tables. The views therefore use
UML conventions where they match the code and identify the persistent entities
as a logical file schema.

\section{Use cases}

Figure~\ref{fig:use-case} records the actions available to the single local
wardrobe user and the external weather dependency. Saving and wardrobe
management remain local actions; the weather service receives a city query but
does not receive wardrobe images.

\begin{figure}[htbp]
  \centering
  \resizebox{0.96\textwidth}{!}{%
  \begin{tikzpicture}[
    font=\small,
    actor/.style={draw=wmgdarkblue!70, rounded corners=2pt, align=center,
      text width=24mm, minimum height=12mm, fill=gray!6},
    usecase/.style={draw=wmgblue!75, ellipse, align=center,
      text width=30mm, minimum height=12mm, fill=wmgblue!7, font=\scriptsize},
    external/.style={draw=green!45!black, rounded corners=2pt, align=center,
      text width=28mm, minimum height=12mm, fill=green!7},
    association/.style={thick, draw=wmgdarkblue!70},
    include/.style={-{Latex[length=2mm]}, dashed, draw=wmgdarkblue!65}
  ]
    \node[actor] (user) at (0,1.0) {Wardrobe user};
    \node[usecase] (city) at (4.1,5.6) {Set or restore city};
    \node[usecase] (select) at (4.1,3.1) {Upload and circle a garment};
    \node[usecase] (analyse) at (9.0,4.1) {Analyse type, colour and style};
    \node[usecase] (context) at (4.1,0.6) {Choose clothing range and style};
    \node[usecase] (recommend) at (9.0,1.6) {Receive complete-outfit recommendation};
    \node[usecase] (save) at (4.1,-1.9) {Save a piece or outfit};
    \node[usecase] (alternatives) at (9.0,-0.9) {Inspect outfit and position alternatives};
    \node[usecase] (manage) at (9.0,-3.4) {Manage local wardrobe};
    \node[external] (weather) at (14.0,5.6) {Open-Meteo geocoding and forecast service};

    \node[draw=wmgdarkblue!45, rounded corners=3pt,
      fit=(city)(select)(analyse)(context)(recommend)(alternatives)(save)(manage),
      inner sep=6mm, label={[font=\bfseries\color{wmgdarkblue}]above:Cove system boundary}] {};

    \draw[association] (user) -- (city);
    \draw[association] (user) -- (select);
    \draw[association] (user) -- (context);
    \draw[association] (user) -- (recommend);
    \draw[association] (user) -- (save);
    \draw[association] (user) -- (manage);
    \draw[association] (weather) -- (city);
    \draw[include] (select) -- node[above,font=\scriptsize]{\guillemotleft include\guillemotright} (analyse);
    \draw[include] (recommend) -- node[right,font=\scriptsize]{\guillemotleft include\guillemotright} (analyse);
    \draw[include] (alternatives) -- node[right,font=\scriptsize]{\guillemotleft extend\guillemotright} (recommend);
  \end{tikzpicture}}
  \caption{UML use-case view of the implemented Cove prototype}
  \label{fig:use-case}
\end{figure}


\section{Overall software structure}

The application is predominantly function-oriented, so an overall class
diagram would invent classes that do not exist. Figure~\ref{fig:overall-architecture}
uses a UML-style component view to show the real source modules, external
service and stored artefacts.

\begin{figure}[htbp]
  \centering
  \resizebox{0.98\textwidth}{!}{%
  \begin{tikzpicture}[
    font=\small,
    component/.style={draw=wmgblue!70, rounded corners=2pt, align=center,
      text width=29mm, minimum height=12mm, fill=wmgblue!7, font=\scriptsize},
    core/.style={component, fill=wmgblue!15, font=\small\bfseries},
    store/.style={component, fill=green!7, draw=green!45!black},
    external/.style={component, fill=gray!8, draw=gray!65},
    arrow/.style={-{Latex[length=2mm]}, thick, draw=wmgdarkblue!70}
  ]
    \node[core, text width=45mm] (app) at (7.2,6.0) {Streamlit presentation and session-state orchestration\\{\ttfamily app.py}};

    \node[component] (recognition) at (0,3.7) {Garment recognition\\{\ttfamily recognition.py}};
    \node[component] (colour) at (3.6,3.7) {Colour analysis\\{\ttfamily colour\_detection.py}};
    \node[component] (weather) at (7.2,3.7) {Weather adapter\\{\ttfamily weather.py}};
    \node[component] (recommendation) at (10.8,3.7) {Recommendation orchestration\\{\ttfamily recommendation.py}};
    \node[component] (wardrobe) at (14.4,3.7) {Wardrobe service\\{\ttfamily wardrobe.py}};

    \node[component] (trained) at (9.0,1.3) {Trained-model adapter\\{\ttfamily trained\_}\\{\ttfamily recommender.py}};
    \node[component] (fallback) at (12.6,1.3) {Deterministic fallback\\Lightweight outfit ranker};
    \node[component] (ranker) at (5.4,1.3) {Candidate search and ranking\\Outfit candidate ranker};

    \node[store] (modelstore) at (4.8,-1.1) {Model checkpoint and abstract prototype catalogue};
    \node[store] (userstore) at (14.4,-1.1) {Local JSON documents and wardrobe images};
    \node[external] (openmeteo) at (8.0,-1.1) {External Open-Meteo APIs};

    \draw[arrow] (app) -- (recognition);
    \draw[arrow] (app) -- (colour);
    \draw[arrow] (app) -- (weather);
    \draw[arrow] (app) -- (recommendation);
    \draw[arrow] (app) -- (wardrobe);
    \draw[arrow] (recommendation) -- (trained);
    \draw[arrow] (recommendation) -- (fallback);
    \draw[arrow] (trained) -- (ranker);
    \draw[arrow] (ranker) -- (modelstore);
    \draw[arrow] (weather) -- (openmeteo);
    \draw[arrow] (wardrobe) -- (userstore);
  \end{tikzpicture}}
  \caption{Overall UML-style component view of the Cove software}
  \label{fig:overall-architecture}
\end{figure}


\section{Recommendation-module classes and dependencies}

Figure~\ref{fig:recommendation-classes} focuses on the classes that are present
in the recommendation path and includes the function-based adapter and facade
on which they depend.

\begin{figure}[htbp]
  \centering
  \resizebox{0.98\textwidth}{!}{%
  \begin{tikzpicture}[
    font=\scriptsize,
    class/.style={draw=wmgdarkblue!65, rounded corners=1pt, align=left,
      text width=42mm, fill=wmgblue!6, inner sep=2.5mm},
    relation/.style={-{Latex[length=2mm]}, thick, draw=wmgdarkblue!70},
    dependency/.style={-{Latex[length=2mm]}, dashed, draw=wmgdarkblue!65}
  ]
    \node[class] (config) at (0,5.0) {\textbf{ModelConfig}\\\hrulefill\\
      embedding\_dim: int = 768\\hidden\_dim: int = 256\\slot\_dim: int = 32\\number\_of\_slots: int = 4\\dropout: float = 0.20};
    \node[class] (model) at (5.3,5.0) {\textbf{CompatibilityRanker}\\\hrulefill\\
      image\_projection: Sequential\\slot\_embedding: Embedding\\scorer: Sequential\\calibration\_temperature: float\\\hrulefill\\
      forward(embeddings, mask)\\score\_preprojected(\\\quad items, mask)\\probability(embeddings, mask)};
    \node[class] (candidate) at (10.6,5.0) {\textbf{OutfitCandidateRanker}\\\hrulefill\\
      device\\model: CompatibilityRanker\\metadata\\last\_search\_stats\\\hrulefill\\
      rank\_exact(fixed, candidates,\\\quad top\_k)\\rank(fixed, candidates, limit)\\\_iter\_combinations(...)};

    \node[class] (adapter) at (5.3,-0.8) {\textbf{Trained recommender adapter}\\\hrulefill\\
      artifacts\_available(...)\\recommend\_with\_trained\_\\\quad model(...)\\weather/style/audience filtering\\abstract-catalogue loading};
    \node[class] (facade) at (10.6,-0.8) {\textbf{Recommendation facade}\\\hrulefill\\
      occupied\_slots\_for\_\\\quad category(...)\\recommend\_outfit(...)\\\_build\_outfit(...)};
    \node[class] (fallback) at (15.9,-0.8) {\textbf{LightweightOutfitRanker}\\\hrulefill\\
      FEATURE\_WEIGHTS\\\hrulefill\\
      features(input, style, slots, colours, weather)\\score(...)\\rank(...)};

    \draw[relation] (config) -- (model);
    \draw[relation] (model) -- (candidate);
    \draw[dependency] (adapter) -- (candidate);
    \draw[dependency] (facade) -- (adapter);
    \draw[dependency] (facade) -- (fallback);
  \end{tikzpicture}}
  \caption{Module-level UML class and dependency diagram for recommendation}
  \label{fig:recommendation-classes}
\end{figure}


\section{Logical entity relationships}

Figure~\ref{fig:persistence-erd} expresses the multiplicities in the local
wardrobe schema. It is an entity-relationship view of serialised application
state, not evidence of a relational database implementation.

\begin{figure}[htbp]
  \centering
  \resizebox{0.96\textwidth}{!}{%
  \begin{tikzpicture}[
    font=\scriptsize,
    entity/.style={draw=wmgdarkblue!65, rounded corners=1pt, align=left,
      text width=34mm, fill=wmgblue!6, inner sep=2.5mm},
    store/.style={entity, fill=green!7, draw=green!45!black},
    rel/.style={thick, draw=wmgdarkblue!70}
  ]
    \node[store] (wardrobe) at (7.0,6.0) {\textbf{WardrobeStore}\\\hrulefill\\version\\pieces[]\\outfits[]\\outfit\_groups[]};
    \node[entity] (piece) at (0,2.7) {\textbf{Piece}\\\hrulefill\\id\\slot\\type and colour\\image\_path\\display\_mode\\source and created\_at};
    \node[entity] (outfit) at (5.0,2.7) {\textbf{Outfit}\\\hrulefill\\id and title\\mood and style\\group\\slots\\created\_at};
    \node[entity] (group) at (10.0,2.7) {\textbf{OutfitGroup}\\\hrulefill\\name\\default\_or\_custom};
    \node[store] (preferences) at (15.0,2.7) {\textbf{UserPreferences}\\\hrulefill\\saved\_cities[0..6]};
    \node[entity] (image) at (0,-0.8) {\textbf{ImageFile}\\\hrulefill\\relative path\\PNG content};
    \node[entity] (slot) at (5.0,-0.8) {\textbf{SlotEntry}\\\hrulefill\\position: one of four slots\\value: embedded Piece or null};

    \draw[rel] (wardrobe) -- node[pos=0.15,left,font=\scriptsize]{1} node[pos=0.85,above,font=\scriptsize]{0..*} (piece);
    \draw[rel] (wardrobe) -- node[pos=0.15,left,font=\scriptsize]{1} node[pos=0.85,right,font=\scriptsize]{0..*} (outfit);
    \draw[rel] (wardrobe) -- node[pos=0.15,right,font=\scriptsize]{1} node[pos=0.85,above,font=\scriptsize]{4..*} (group);
    \draw[rel] (outfit) -- node[pos=0.15,left,font=\scriptsize]{1} node[pos=0.85,right,font=\scriptsize]{4} (slot);
    \draw[rel] (slot) -- node[above,font=\scriptsize]{0..1 snapshot} (piece);
    \draw[rel] (piece) -- node[right,font=\scriptsize]{0..1} (image);
    \draw[rel] (group) -- node[pos=0.15,above,font=\scriptsize]{1} node[pos=0.85,above,font=\scriptsize]{0..*} (outfit);
  \end{tikzpicture}}
  \caption{Entity-relationship view of the logical file schema; these entities are serialised to JSON rather than relational tables}
  \label{fig:persistence-erd}
\end{figure}


\section{User stories and acceptance criteria}

Table~\ref{tab:user-stories} converts the implemented interaction paths into
testable user stories. Priorities distinguish the core recommendation journey
from supporting convenience features.

\begingroup
\scriptsize
\setlength{\tabcolsep}{3pt}
\renewcommand{\arraystretch}{1.18}
\begin{longtable}{@{}p{0.06\textwidth}p{0.13\textwidth}p{0.29\textwidth}p{0.38\textwidth}p{0.07\textwidth}@{}}
  \caption{Implemented user stories and acceptance criteria}\label{tab:user-stories}\\
  \toprule
  \textbf{ID} & \textbf{Feature} & \textbf{User story} & \textbf{Acceptance criteria} & \textbf{Priority} \\
  \midrule
  \endfirsthead
  \multicolumn{5}{l}{\textit{Table~\thetable\ continued}}\\
  \toprule
  \textbf{ID} & \textbf{Feature} & \textbf{User story} & \textbf{Acceptance criteria} & \textbf{Priority} \\
  \midrule
  \endhead
  \bottomrule
  \endfoot
  US-01 & Location & As a user, city weather should inform the outfit context. & A valid city returns current conditions; an invalid or unavailable request produces a bounded error rather than stopping the application. & Must \\
  US-02 & Saved cities & As a user, frequently used cities should be available without retyping. & A city can be saved or removed; no more than six are retained; selecting one refreshes the forecast. & Should \\
  US-03 & Image selection & As a user, a garment should be selectable from a photograph. & The user can upload an image, draw a region and repeat the selection before analysis. & Must \\
  US-04 & Recognition & As a user, the selected garment should be interpreted before recommendation. & The application returns a supported broad category, colour, style cues and an image embedding, with an editable fallback when recognition is uncertain. & Must \\
  US-05 & Context control & As a user, clothing range and style should remain under explicit control. & Menswear or womenswear and an available style can be selected; changing either regenerates the recommendation without inferring gender from the image. & Must \\
  US-06 & Outfit completion & As a user, the selected garment should remain fixed while the outfit is completed. & The result covers all feasible supported positions, respects occupied slots and weather constraints, and uses product-independent types and colours. & Must \\
  US-07 & Alternatives & As a user, more than one acceptable solution should be inspectable. & The interface provides a primary outfit, a materially different alternative and model-ranked position replacements where feasible. & Must \\
  US-08 & Operational fallback & As a user, recommendation should remain available if model artefacts cannot load. & The system records the loading reason and returns a deterministic rule-ranked result without presenting it as a learned score. & Must \\
  US-09 & Save result & As a user, either the supplied garment or a complete recommendation should be saved. & Saving preserves the user image where available and represents abstract recommended garments as coloured icons. & Should \\
  US-10 & Outfit editing & As a user, saved outfits should be completed and revised later. & Each of four equal positions supports add, replace or confirmed removal; an outfit can be renamed and assigned to a group. & Should \\
  US-11 & Piece library & As a user, individual garments should be reusable across outfits. & Pieces can be filtered by four positions, renamed, reused in multiple outfits and deleted with confirmation. & Should \\
  US-12 & Local persistence & As a user, wardrobe information should remain private on the device. & Wardrobe JSON, preferences and images are stored locally and are not incorporated into model training. & Must \\
\end{longtable}
\endgroup
